\section{Applicability, Limitations, and Future Extensions}

\renewcommand{\thefigure}{F.\arabic{figure}}  % Number figures as A.1, A.2, ...
\renewcommand{\thetable}{f.\arabic{table}}  % Number tables as A.1, A.2, ...


Sensitivity-constrained training requires Jacobians. This can be challenging in systems where adjoints are difficult to implement, unstable, infeasible, or prohibitively expensive, for example, highly chaotic or turbulent dynamics, nonlinear elastodynamics, or other mathematical non-differentiabilities. In such settings, operational systems may rely on large-scale ensembles or extreme-scale offline computation to retain full physical fidelity. Nevertheless, even in chaotic and turbulent systems, adjoint sensitivities are used operationally over the finite time windows in which they remain stable, as in 4D-Var and ensemble-variational data assimilation. SC-NO can therefore be applicable to a broad class of PDE-governed systems with locally stable adjoints.

While our SC-NO tsunami workflow is robust in the face of inversion uncertainty and moderate buoy data noise, it is not intended as an operational forecasting system, which must also account for additional uncertainties in rupture geometry, background currents, non-seismic sources, sensor outages, and near-coast inundation physics \citep{grezio2017probabilistic}. Future work can also incorporate seismic and sea floor sensor data to improve lead time. Rather than proposing a full operational paradigm, our current aim is to isolate a complementary computational contribution: real-time surrogate evaluation and minute-scale inversion. The method can, in principle, enable faster exploration of such uncertainties and facilitate data assimilation, particularly in settings where repeated full-physics solves would be prohibitive, although it does not replace existing operational pipelines. 

Prior studies have shown that a method that incorporates PDE equation loss like PINO can provide moderate improvements for state accuracy but does not directly constrain parameter sensitivities and introduces additional computational cost \citep{behroozi2025sensitivity} and future work can test its cost and benefit. Because our focus is on full-field, high-resolution surrogates that remain accurate under localized input perturbations and enable stable inversion, we did not compare against traditional reduced-order approaches such as Proper Orthogonal Decomposition, which are most effective when low-dimensional manifolds exist, nor against physics-informed neural networks (PINNs), which target different objectives such as knowledge discovery. Dimensionality-reduction and gradient-aware methods such as DINO or DE-DeepONet address complementary challenges—primarily representational efficiency—whereas SC-NO targets sensitivity, fidelity, and inverse stability. These approaches are orthogonal in design and span different axes of the accuracy–cost–conditioning trade space, and assessing whether low-dimensional representations preserve the reverse-scaling behavior identified here remains an important direction for future work.

Future directions include adaptive constraint weighting, higher-order (Hessian) supervision, uncertainty quantification via learned Jacobians, and adaptive search for optimal Jacobian sampling strategies. Collectively, these results position sensitivity-constrained training as a foundational enhancement to neural operators, enabling physically trustworthy surrogates for prediction, inversion, and decision-making under uncertainty.
